% ============ 第2章 §2.19 Carleson插值问题 (II)(原书 pp.185-190) ============
\section{Carleson 插值问题 (II)}\label{sec:2.19}

本节证明命题~\ref{prop:2.18.4} 的充分性.此时由引理~\ref{lem:2.18.6} 完成命题~\ref{prop:2.18.4} 的证明.另外给出一致分离的 $\{z_{j}\}_{j=1}^{\infty}$ 的例子.

\begin{lemma}\label{lem:2.19.1}
若 $\rho_{j} > 0$($j \geq 1$),则
\begin{equation*}
\sum_{j=1}^{\infty}\frac{(1 - |z_{j}|^{2})(1 - |z_{k}|^{2})}{|1 - \bar{z}_{j}z_{k}|^{2}} \leq 1 - 2\log\rho_{k} \quad (k \geq 1).
\end{equation*}
\end{lemma}

\begin{proof}
令 $\alpha_{jk} = (1 - |z_{j}|^{2})(1 - |z_{k}|^{2})/|1 - \bar{z}_{j}z_{k}|^{2}$,由简单计算,
\begin{equation*}
\left| \frac{z_{k} - z_{j}}{1 - \bar{z}_{j}z_{k}} \right|^{2} = 1 - \alpha_{jk}.
\end{equation*}
于是
\begin{align*}
\rho_{k}^{2} &= \prod_{j \neq k}\left| \frac{z_{k} - z_{j}}{1 - \bar{z}_{j}z_{k}} \right|^{2} = \prod_{j \neq k}(1 - \alpha_{jk}) \\
&\leq \exp\left\{-\sum_{j \neq k}\alpha_{jk}\right\},
\end{align*}
利用 $0 < \alpha_{jk} \leq 1$ 即证引理.
\end{proof}

\begin{lemma}\label{lem:2.19.2}
设 $a_{jk} \in \mathbb{C}$($1 \leq j, k \leq n$)满足 $a_{kj} = \overline{a_{jk}}$ 且
\begin{equation*}
\sum_{j=1}^{n}|a_{jk}| \leq M \quad (1 \leq k \leq n).
\end{equation*}
此时对任意 $\{x_{j}\}_{j=1}^{n} \subset \mathbb{C}$,
\begin{equation*}
\left| \sum_{j,k=1}^{n} a_{jk}x_{j}\bar{x}_{k} \right| \leq M\sum_{j=1}^{n}|x_{j}|^{2}.
\end{equation*}
\end{lemma}

\begin{proof}
用 Schwartz 不等式,
\begin{align*}
\sum_{j,k=1}^{n}|a_{jk}x_{j}\bar{x}_{k}|&\leq \sum_{k=1}^{n}\left(\sum_{j=1}^{n}|a_{jk}||x_{j}|^{2}\right)^{1/2}\left(\sum_{j=1}^{n}|a_{jk}||x_{k}|^{2}\right)^{1/2}  \\
&\leq M^{1/2}\sum_{k=1}^{n}\left(\sum_{j=1}^{n}|a_{jk}||x_{j}|^{2}\right)^{1/2}|x_{k}| \\
&\leq M^{1/2}\left\{\sum_{k=1}^{n}\left(\sum_{j=1}^{n}|a_{jk}||x_{j}|^{2}\right)\right\}^{1/2}\left(\sum_{k=1}^{n}|x_{k}|^{2}\right)^{1/2} \\
&= M^{1/2}\left\{\sum_{j=1}^{n}\left(\sum_{k=1}^{n}|a_{jk}|\right)|x_{j}|^{2}\right\}^{1/2}\left(\sum_{k=1}^{n}|x_{k}|^{2}\right)^{1/2} \\
&\leq M\left(\sum_{j=1}^{n}|x_{j}|^{2}\right)^{1/2}\left(\sum_{k=1}^{n}|x_{k}|^{2}\right)^{1/2}.
\end{align*}
\end{proof}

如 2.17 节那样在 $H^{2}$ 上定义内积 $\langle\ ,\ \rangle$.即对 $f, g \in H^{2}$,
\begin{equation*}
\langle f, g \rangle = \frac{1}{2\pi}\int_{-\pi}^{\pi} f(e^{i\theta})\overline{g(e^{i\theta})}\dd \theta.
\end{equation*}
设 $B_{n}(z)$ 与 $B_{nk}(z)$($k \leq n$)为 2.17 节定义的 Blaschke 积.证明当 $\inf_{j}\rho_{j} > 0$ 时 $\{((1 - |z_{j}|^{2})^{1/2}f(z_{j})); f \in H^{2}\} = \ell^{2}$.

设 $w = (w_{k}) \in \ell^{2}$ 为任意元素,令
\begin{equation*}
F_{nk}(z) = (1 - |z_{k}|^{2})^{3/2}[B_{n}(z)]^{2}(z - z_{k})^{-2}
\end{equation*}
且
\begin{equation*}
f_{n}(z) = \sum_{k=1}^{n} w_{k}[B_{nk}(z_{k})]^{-2}F_{nk}(z).
\end{equation*}
此时
\begin{equation*}
(1 - |z_{k}|^{2})^{1/2}f_{n}(z_{k}) = w_{k} \quad (1 \leq k \leq n)
\end{equation*}
且
\begin{align*}
& \|f_{n}\|_{2}^{2} = \langle f_{n}, f_{n} \rangle \\
& \qquad = \sum_{j,k=1}^{n} w_{j}[B_{nj}(z_{j})]^{-2}\bar{w}_{k}[\overline{B_{nk}}(z_{k})]^{-2}\langle F_{nj}, F_{nk} \rangle .
\end{align*}
由留数定理,
\begin{equation*}
\langle F_{nj}, F_{nk} \rangle = (1 - |z_{j}|^{2})^{3/2}(1 - |z_{k}|^{2})^{3/2}(1 + z_{j}\bar{z}_{k})(1 - z_{j}\bar{z}_{k})^{-3},
\end{equation*}
由 $(1 - |z_{j}|^{2})(1 - |z_{k}|^{2}) \leq |1 - z_{j}\bar{z}_{k}|^{2}$,
\begin{equation*}
|\langle F_{nj}, F_{nk} \rangle| \leq 2(1 - |z_{j}|^{2})(1 - |z_{k}|^{2})|1 - z_{j}\bar{z}_{k}|^{-2}.
\end{equation*}
由引理~\ref{lem:2.19.1},
\begin{equation*}
\sum_{j=1}^{n}|\langle F_{nj}, F_{nk} \rangle| \leq 2(1 - 2\log\rho_{k}) \quad (1 \leq k \leq n),
\end{equation*}
由引理~\ref{lem:2.19.2},
\begin{align*}
\|f_{n}\|_{2}^{2} &\leq \sum_{j,k=1}^{n}|w_{j}|\rho_{j,n}^{-2}|w_{k}|\rho_{k,n}^{-2}|\langle F_{nj}, F_{nk} \rangle| \\
&\leq 2(1 - 2\log\inf_{1 \leq k \leq n}\rho_{k})\sum_{j=1}^{n}|w_{j}|^{2}\rho_{j,n}^{-4} \\
&\leq 2(1 - 2\log\inf_{1 \leq k \leq n}\rho_{k})\sup_{1 \leq j \leq n}\rho_{j}^{-4}\sum_{j=1}^{n}|w_{j}|^{2} \\
&\leq 2(1 - 2\log\inf_{j}\rho_{j})(\inf_{j}\rho_{j})^{-4}\sum_{j=1}^{\infty}|w_{j}|^{2}.
\end{align*}
如此由引理~\ref{lem:2.10.1} 可证 $\{f_{n}\}$ 是正规族.于是存在 $f \in H^{2}$,对每个 R < 1,$\{f_{n}\}$ 的某个子列在 $\{z; |z| < R\}$ 上一致收敛.于是
\begin{equation*}
(1 - |z_{j}|^{2})^{1/2}f(z_{j}) = w_{j} \quad (j \geq 1)
\end{equation*}
且
\begin{equation*}
\|f\|_{2} \leq 2(1 - 2\log\inf_{j}\rho_{j})(\inf_{j}\rho_{j})^{-2}\|w\|_{2}.
\end{equation*}
如此证明了 $\{((1 - |z_{j}|^{2})^{1/2}f(z_{j})); f \in H^{2}\} \supset \ell^{2}$.

还必须证明 $\{((1 - |z_{j}|^{2})^{1/2}f(z_{j})); f \in H^{2}\} \subset \ell^{2}$.对任意 $w = (w_{j}) \in \ell^{2}$,由前半的证明,设 $\gamma = 2(1 - 2\log\inf_{j}\rho_{j})(\inf_{j}\rho_{j})^{-2}$,则
\begin{equation*}
\min \{\|g\|_{2}; g \in H^{2}, (1 - |z_{j}|^{2})^{1/2}g(z_{j}) = w_{j} \ (1 \leq j \leq n)\} \leq \gamma\|w\|_{2}.
\end{equation*}
于是由定理~\ref{thm:2.17.1},任意 $f \in H^{2}$ 满足下面的不等式:
\begin{equation*}
\left| \sum_{j=1}^{n} \frac{w_{j}(1 - |z_{j}|^{2})^{1/2}}{B_{nj}(z_{j})} f(z_{j}) \right| \leq \gamma\|w\|_{2}\|f\|_{2}.
\end{equation*}
这里利用 $(\ell^{2})^{*} = \ell^{2}$,对任意 n,
\begin{equation*}
\left(\sum_{j=1}^{n}\frac{(1 - |z_{j}|^{2})}{\rho_{j,n}^{2}}|f(z_{j})|^{2}\right)^{1/2} \leq \gamma\|f\|_{2}.
\end{equation*}
于是设 $\gamma_{0} = 2(1 - 2\log\inf_{j}\rho_{j})(\inf_{j}\rho_{j})^{-1}$,则
\begin{equation*}
\left(\sum_{j=1}^{\infty}(1 - |z_{j}|^{2})|f(z_{j})|^{2}\right)^{1/2} \leq \gamma_{0}\|f\|_{2},
\end{equation*}
故导出 $\{((1 - |z_{j}|^{2})^{1/2}f(z_{j})); f \in H^{2}\} \subset \ell^{2}$.于是与前半合在一起 $\{((1 - |z_{j}|^{2})^{1/2}f(z_{j})); f \in H^{2}\} = \ell^{2}$.

接下来的 Hayman \cite{ch2:19} 与 Newman \cite{ch2:20} 的例~\ref{exa:2.19.4} 回答了是否存在一致分离的 $\{z_{j}\}_{j=1}^{\infty}$ 这一问题.

\begin{lemma}\label{lem:2.19.3}
若 $|\alpha| < 1$ 且 $|\beta| < 1$,则
\begin{equation*}
\left| \frac{\alpha - \beta}{1 - \bar{\alpha}\beta} \right| \geq \frac{|\alpha| - |\beta|}{1 - |\alpha\beta|}.
\end{equation*}
\end{lemma}

\begin{proof}
当 $1 \leq A < B$ 时,
\begin{equation*}
\frac{A + x}{B + x} \geq \frac{A - 1}{B - 1} \quad (-1 \leq x \leq 1),
\end{equation*}
令 a > 0,$A = \frac{a^{2} + r^{2}}{2ar}$,$B = \frac{1 + a^{2}r^{2}}{2ar}$ 且 $x = -\cos\theta$,则
\begin{equation*}
\left| \frac{a - re^{i\theta}}{1 - are^{i\theta}} \right|^{2} = \frac{A + x}{B + x} \geq \frac{A - 1}{B - 1} = \left(\frac{a - r}{1 - ar}\right)^{2}.
\end{equation*}
当 $a = \alpha$,$re^{i\theta} = \beta$ 时即证引理,一般 $\alpha$ 的情形容易导出.
\end{proof}

\begin{example}\label{exa:2.19.4}
当 $\{z_{j}\}_{j=1}^{\infty} \subset D$ 时,若存在 0 < c < 1 使得
\begin{equation*}
1 - |z_{n}| \leq c(1 - |z_{n-1}|) \quad (1 \leq n < \infty),
\end{equation*}
则 $\{z_{j}\}_{j=1}^{\infty}$ 是一致分离的.
\end{example}

\begin{proof}
当 j > k 时,由假设,
\begin{equation*}
1 - |z_{j}| \leq c^{j-k}(1 - |z_{k}|) \quad (1 \leq k < \infty).
\end{equation*}
于是
\begin{equation*}
|z_{j}| - |z_{k}| \geq (1 - c^{j-k})(1 - |z_{k}|)
\end{equation*}
且
\begin{align*}
1 - |z_{j}z_{k}| & = 1 - |z_{j}| + |z_{j}|(1 - |z_{k}|) \\
& \leq (1 + c^{j-k})(1 - |z_{k}|).
\end{align*}
因此由引理~\ref{lem:2.19.3},当 j > k 时,
\begin{equation*}
\left| \frac{z_{k} - z_{j}}{1 - \bar{z}_{j}z_{k}} \right| \geq \frac{|z_{j}| - |z_{k}|}{1 - |z_{j}z_{k}|} \geq \frac{1 - c^{j-k}}{1 + c^{j-k}}.
\end{equation*}
当 j < k 时,上面的不等式为
\begin{equation*}
\left| \frac{z_{k} - z_{j}}{1 - \bar{z}_{j}z_{k}} \right| \geq \frac{1 - c^{k-j}}{1 + c^{k-j}}.
\end{equation*}
如此
\begin{align*}
\prod_{\substack{j=1\\ j\neq k}}^{\infty}\left| \frac{z_{k} - z_{j}}{1 - \bar{z}_{j}z_{k}} \right| &= \prod_{j=1}^{k-1}\left| \frac{z_{k} - z_{j}}{1 - \bar{z}_{j}z_{k}} \right| \times \prod_{j=k+1}^{\infty}\left| \frac{z_{k} - z_{j}}{1 - \bar{z}_{j}z_{k}} \right| \\
&\geq \prod_{n=1}^{k-1}\frac{1 - c^{n}}{1 + c^{n}}\prod_{n=1}^{\infty}\frac{1 - c^{n}}{1 + c^{n}}.
\end{align*}
因 $\sum_{n=1}^{\infty}2c^{n}/(1 + c^{n}) < \infty$,上面的无限乘积 $\prod_{n=1}^{\infty}(1 - c^{n})/(1 + c^{n}) > 0$,故 $\{z_{j}\}_{j=1}^{\infty}$ 是一致分离的.
\end{proof}

Carleson 插值问题针对的是正则函数,关于调和函数的研究见 \cite{ch2:21}.
